% ============================================================================
% PART I: CONSTRUCTION
% Section 4: TQFT and Its Observables
% ============================================================================
%
% Purpose
%   After Secs. 2--3 the reader already knows enough about CS theory to see
%   that the bordism functor story is not an empty abstraction. Use this
%   section to (a) state the axiomatic / functorial definition, (b) briefly
%   distinguish classes of TQFTs (Schwarz vs. Witten, abelian vs. nonabelian),
%   (c) identify the canonical observables --- Wilson lines, linking/framing,
%   Hilbert spaces on surfaces, ground-state degeneracy, knot invariants ---
%   and (d) record the observable-family taxonomy used in Part II.
%
% Status: MERGE from archive + write new:
%   - NEW: Atiyah-Segal functorial summary (brief, ~3--4 pages).
%   - EXTRACT from archive/chern_simons_two_review_synthesis.tex
%     for Wilson-loop and canonical-quantization narrative.
%   - EXTRACT from archive/dense_derivation_expansion.tex
%     for abelian linking, torus H(T^2), framing calculation.
%   - EXTRACT from writings/brian_chern_simons_theory.tex
%     for symplectic / geometric quantization rigor.
%
% Length target: ~15--20 pages
%
% Subsections
%   4.1  The functorial definition (Atiyah--Segal)
%   4.2  Classes of TQFTs (Schwarz-type vs Witten-type, 2d, 3d, 4d examples)
%   4.3  Chern-Simons as the central example
%   4.4  Wilson lines, linking, and framing
%   4.5  Hilbert spaces on surfaces and ground-state degeneracy
%   4.6  Observable families: taxonomy for Part II
%
% References
%   Atiyah 1988 (\cite{Atiyah1988TQFT})                 -- axioms
%   Witten 1988 (\cite{Witten1988TQFT})                 -- Witten-type
%   Witten 1989 (\cite{Witten1989Jones})                -- CS, Jones, framing
%   Blau-Thompson 1991 (\cite{BlauThompson1991TopologicalGauge}) -- BF, beyond CS
%   Birmingham et al. 1991 (\cite{BirminghamBlauRakowskiThompson1991TopologicalFieldTheory})
%   Dunne 1999 (\cite{Dunne1999AspectsCS})              -- physics review
%   Dijkgraaf 1989 (\cite{Dijkgraaf1989Geometry2dCFT})  -- 2d/3d boundary
%   Elitzur-Moore-Schwimmer-Seiberg 1989 (\cite{ElitzurMooreSchwimmerSeiberg1989CanonicalCS})
%                                                       -- canonical quantization
%   Wen-Niu 1990 (\cite{WenNiu1990GroundStateDegeneracy}) -- GSD on high genus
%
% Writing notes
%   - The section is the last step of Part I. Close with a short paragraph
%     introducing the observable taxonomy (transport / spectroscopic /
%     braiding / defect) that Part II will organize around.

\section{Topological Quantum Field Theory and Its Observables}
\label{sec:tqft-observables}

% TODO: Preamble framing --- "the functorial definition codifies what the
% Chern-Simons path integral is already doing."

\subsection{The Functorial Definition}
\label{subsec:functorial-definition}

% TODO: Bord_n as symmetric monoidal category. Z: Bord_n -> Vect_k
% assigns vector spaces to closed (n-1)-manifolds and linear maps to bordisms.
% Gluing = composition; disjoint union = tensor product.
% Cite: \cite{Atiyah1988TQFT},
%       \cite{BirminghamBlauRakowskiThompson1991TopologicalFieldTheory}

\subsection{Classes of TQFTs}
\label{subsec:tqft-classes}

% TODO: Schwarz-type (CS, BF) vs Witten-type (cohomological);
% 2d TQFTs and Frobenius algebras (brief); 3d CS as main example;
% a word about 4d Donaldson/Seiberg-Witten for completeness (no derivation).
% Cite: \cite{Witten1988TQFT}, \cite{BlauThompson1991TopologicalGauge},
%       \cite{BirminghamBlauRakowskiThompson1991TopologicalFieldTheory}

\subsection{Chern--Simons as the Central Example}
\label{subsec:cs-central-example}

% TODO: S^3 partition function, genus-g Hilbert spaces, Wilson-line correlators
% as knot invariants. Abelian U(1)_k: explicit and Gaussian. Nonabelian G_k:
% reference to quantum groups, Verlinde formula.
% Cite: \cite{Witten1989Jones}, \cite{Dunne1999AspectsCS}

\subsection{Wilson Lines, Linking, and Framing}
\label{subsec:wilson-linking-framing}

% TODO: W_R(K) = tr_R P exp(oint_K A). Abelian: <W_K W_L> = exp(2 pi i Lk(K,L)/k).
% Framing dependence as genuine quantum data, not a bug.
% Source: archive/dense_derivation_expansion.tex worked calculation of abelian
% linking.
% Cite: \cite{Witten1989Jones}, \cite{Dunne1999AspectsCS}

\subsection{Hilbert Spaces on Surfaces and Ground-State Degeneracy}
\label{subsec:hilbert-surfaces-gsd}

% TODO: Phase space on Sigma = M_flat(Sigma, G), quantization -> H(Sigma).
% Torus example: dim H_U(1)_k(T^2) = k (worked calculation).
% High-genus generalization dim H = k^g (abelian).
% Source: archive/dense_derivation_expansion.tex, writings/brian_chern_simons_theory.tex.
% Cite: \cite{ElitzurMooreSchwimmerSeiberg1989CanonicalCS},
%       \cite{WenNiu1990GroundStateDegeneracy}

\subsection{An Observable Taxonomy}
\label{subsec:observable-taxonomy}

% TODO: Introduce the four observable families used to organize Part II.
% From tqft_observables_literature_review.md Sec. 2:
%   (i)  Transport / response (Hall conductance, chiral central charge,
%        magnetoelectric response).
%   (ii) Spectroscopic / CP-sensitive (topological susceptibility, eta' mass,
%        vacuum energy vs theta, neutron EDM, axion).
%   (iii) Braiding / nonlocal (Wilson-loop VEVs, linking, framing, anyon
%         exchange).
%   (iv) Defect / global-structure (magnetic charge, Dirac quantization,
%        monopoles, 't Hooft lines, global form of gauge group).
% Cite: \cite{Stern2008AnyonsQHEReview},
%       \cite{Witten2016ThreeLecturesTopologicalPhases}

% Transition: Part II organizes the physics by observable family, not by
% subfield.
